% !TeX root = ../papers/nc-p3.tex
\section{十八次の二次錐の頂点と長さ一の部分正規化}
\label{sec:cofactor-eighteen-cone-inverse}
\begin{lemma}\label{lem:cofactor-colength-one-multiplicity}
被約なGorenstein曲線局所環 $R$ に対し余長一の有限双有理拡大 $R\subset R'$ があり、
$R'$ の各局所環もGorensteinなら、$R$ の重複度は二である（剰余体は代数閉）。
\end{lemma}
\begin{proof}
導手は $R'/R$ の零化イデアルなので $R$ の極大イデアル $\mathfrak m$ である。
有限双対性により $\operatorname{Hom}_R(R',R)=\mathfrak m$ は $R'$ の標準加群である。
従って $\mathfrak m$ は半局所環 $R'$ の可逆イデアルで、非零因子を生成元とする主イデアルとなる。
$R'$ 上の $\mathfrak m$ の重複度は従って $\ell_R(R'/\mathfrak m)=2$ である。
有限余長差は一次元の重複度を変えないので $R$ 上でも二となる。
\end{proof}

\begin{theorem}\label{thm:cofactor-eighteen-cone-not-birational}
$d=18$ で核像が二次錐のとき、核写像は次数二である。
\end{theorem}
\begin{proof}
双有理と仮定する。核切断 $J$ は錐 $Q$ 上のCartier三次切断である。
実際錐の重み $1,1,2$ の座標では整な六次曲線の方程式は重み六で、普通の三次式に持ち上がる。
逆四空間も $\mathcal O_Q(3)$ の切断で実現される。実際四つの入力座標の比に共通固定因子を残さない
階数一反射層を取れば一般の零Weil因子が $J$ で、その類は $3L$ である。
ここで $L$ は核超平面である。$p_a(J)=4$ なので定理\ref{thm:cofactor-sectional-graph-genus}の
$G\to J$ の余長は一、その唯一の非自明な点 $q$ は逆基点である。一般切断上でその重複度は二である。
実際一般入力平面で $q$ 上の入力点は階数三である。$q$ を消す固定列の曲線部分は次数二以下で $C$ と
共通成分を持たず、有限個の基点候補全てについて一般平面はこの有限例外を避ける。
階数三では入力から核点が正則に復元され、$G$ は局所的に入力平面曲線に等しい。
従って補題\ref{lem:cofactor-colength-one-multiplicity}が使える。

$q$ が滑らかな点なら、そこを一度吹いた $S_1=\operatorname{Bl}_q Q$ で
$J'$ の算術種数が一だけ下がる。$G=J'$ はGorenstein曲線の余長一の中間環の一意性による。
残余基点 $Z$ は $S_1$ の特異点でも局所完全交叉である。実際一般切断 $a$ を法として
入力比が正則なので基点イデアルは二生成であり、曲面はCohen--Macaulayである。
$A'=3L-2e$ から長さは $A'^2-8=6$。$\Gamma=\operatorname{Bl}_Z S_1$ 上の入力束を $H$ とする。
$q$ からの射影 $h=L-e$ は平面への双有理射で、潰す曲線は $q$ を通る唯一の錐母線だけである。
$A'-h=2L-e$ の八次元の切断に長さ六の条件を課しても二独立切断が残る。
それらの平面像は三次因子であり、独立な二元の像も独立である（双有理写像の関数体で比を取る）。
従って $h$ で正次数を持つ入力縮約曲線は高々二本である。
実際その曲線では $H-h$ が負次数なので全切断の共通成分となる。
例外直線の入力次数は一で、正次元繊維の入力像は合計高々三点である。

滑らかな二次曲面とその一点および相異なる六残余点のモデルへ変形できることを確認する。
$q$ は滑らかな点なので錐の平滑化とともに動かせる。
特異点に台を持つ部分も含め $Z$ は零次元局所完全交叉であり、相対Hilbertスキームはここで底へ滑らかである。
局所生成正則列は底の微小拡大上に持ち上がり、局所持ち上げの差の障害は
$H^1(Z,\mathcal Hom(I/I^2,\mathcal O_Z))=0$。平坦性は正則列の繊維判定による。
残余のブローアップも一方程式の相対表示で平坦な局所完全交叉である。
$J'=G$ の算術種数三、入力束次数八から $H^1(I_Z A')=0$ を得る
（$H^1(\mathcal O_{S_1})=H^2(\mathcal O_{S_1})=0$ で、一般切断を除いた商は $J'$ 上の線束）。
従って四入力切断を平滑化できる。一般の四切断と中心へ移動して得るのは
定理\ref{thm:cofactor-eighteen-birational-quadric}の平面八点の四次モデルである。
同じEuler数から入力有限な開集合で $M=f_*L$。少なくとも十七個の孤立三点繊維が
一般モデルの長さ十三に入らねばならず矛盾する。

$q$ が頂点の場合、最小解消 $S_0=\mathbb F_2$ で $L=E+2F$ とする。
一般切断の固有変換 $J_0$ は $uE+6F$ で、頂点での $J$ の重複度は $J_0E=6-2u$ である。
実際頂点イデアルはこの解消で $\mathcal O(-E)$ に引き戻され、曲線の正規化上で
一般パラメータの位数を足したものが重複度である。
前半の重複度二より $u=2$。$J_0$ の算術種数は三なのでやはり $J_0=G$、
$S_0$ 上の全基点は長さ八の局所完全交叉となる。
逆束は $2E+6F=2L+2F$、母線への直像は次数 $2,4,6$ の三直和である。
定理\ref{thm:cofactor-eighteen-birational-quadric}の $(2,4)$ の議論がここでも成立する。
各母線との基点スキームの交点長は二以下、長さ八の条件後の直像は非負次数の階数三、全次数四の束 $W$ である。
像は $\mathbb P(W)$ 上の $2\xi$ の因子なので垂直縮約の入力点は高々四、水平縮約は高々二である。
$H^1(I_Z(2L+2F))=0$、$\mathbb F_2$ と $(L,F)$ を積曲面へ動かし、
長さ八の局所完全交叉と入力四切断も一緒に動かせる。従って有限な開集合での同じ押出し同型から来る
少なくとも十四個の孤立三点繊維が長さ十二に入る必要があり矛盾する。
\end{proof}
\path{computations/verify_cofactor_eighteen_cone.py} は頂点重複度、固有変換の種数および二種類の平滑化の交点数を検算する。
